\scanpage{97}
arête basique \inserted{(disons $\unsure{a|u}=\{\alpha,\beta\}$)}, l'application
\[
\widetilde\Gamma\xrightarrow{\widetilde\varphi_u}\prod_{i\in u}\widetilde E_i
\quad(\simeq\widetilde E_\alpha\times\widetilde E_\beta)
\]
est \uline{injective}, et son image est formée des \deleted{\unsure{toutes|tous}} les \deleted{familles} couples $(s,t)\in\widetilde E_\alpha\times\widetilde E_\beta$ tels que $s$ soit lié à $t$.

\uline{Scholie} En résumé, moyennant les conditions \textcircled{1} et \textcircled{2} plus haut, on voit que la structure de graphe $S$, muni du triangle $t_0$, se reconstitue à partir de la donnée de

a) un ens $t_0$ tel que $\operatorname{Card}t_0=3$

b) une famille $\deleted{\unsure{E|\widetilde E}}F=(\widetilde E_\alpha)_{\alpha\in t_0}$ d'ens. \inserted{munis chacun d'une invol. $\sigma_\alpha$ ss pt fixe}

c) une partie $\widetilde\Gamma\lhook\joinrel\xrightarrow[\widetilde\varphi]{}\underbrace{\prod_{\alpha\in t_0}\widetilde E_\alpha}_{\text{\textquotedbl}\widetilde p}$

[satisfaisant les conditions

$\alpha$) $\forall u\in\unsure{\mathcal P|\mathfrak P}_2(t_0)$, l'application $\widetilde\Gamma\xrightarrow{\widetilde\varphi_u}\prod_{\alpha\in u}\widetilde E_\alpha$ est injective,

\inserted{(\unsure{**}) Si $\alpha,\beta\in t_0$, $\alpha\ne\beta$, et si $s\in\widetilde E_\alpha$, $t\in\widetilde E_\beta$, si $(s,t)\in\mathrm{Im}\,\widetilde\varphi_{\{\alpha,\beta\}}(\widetilde\Gamma)$ alors $(\sigma_\alpha s,t)\notin\mathrm{Im}\,\widetilde\varphi_{\{\alpha,\beta\}}(\widetilde\Gamma)$}
% margin note written sideways
\margin{$\boxed{**}$ Si $(s_\alpha)_{\alpha\in\underline t_0}\in\prod_{\alpha\in\underline t_0}\widetilde E_\alpha$ \deleted{est tel que} \deleted{\unsure{$\forall\alpha$ si}} $(\underline t_\alpha)_{\alpha\in\underline t_0}\in{\unsure{\widetilde T|\widetilde\Gamma}}^{\underline t_0}$ sont tels que \deleted{\unsure{$\forall$}} $\alpha,\beta\in\underline t_0$, $\alpha\ne\beta$ on a $s_\beta=\varphi_\beta(\underline t_\alpha)$, alors les $\underline t_\alpha$ sont égaux}

de la façon suivante

\[
S\simeq t_0\amalg\coprod_{\alpha\in t_0}\widetilde E_\alpha \tag*{$1^\circ$)}
\]
\[
2^\circ)\ \deleted{\text{deux pts}}\ \left\{\begin{aligned}
&\alpha)\ \text{Deux pts de }t_0\text{ sont liés}\\
&\beta)\ \text{un point de }\widetilde E_\alpha\text{ est lié à un et un}\\
&\quad\text{seul él. de }t_0,\ \text{savoir }\alpha\quad(\text{si }\alpha\in t_0)\\
&\unsure{\gamma) }\text{un point }s\text{ de }\overline E_\alpha\text{ est lié à un }\inserted{\text{et un}}\text{ seul pt de}\\
&\quad\unsure{\text{lui}|\widetilde E_\alpha},\ \text{savoir }\sigma_\alpha(s)\\
&\delta)\ \deleted{\text{Soit}}\ \alpha,\beta\in t_0,\ \alpha\ne\beta,\quad s\in\widetilde E_\alpha,\ t\in\overline E_\beta,\quad\text{alors }s,t\text{ sont}
\end{aligned}\right.
\]
% archive sheet number: 96
